\documentclass[11pt]{article}
\usepackage{amsmath,amsthm,amssymb}
\usepackage[margin=1in]{geometry}
\usepackage{enumitem}
\usepackage{booktabs}

\newtheorem{theorem}{Theorem}
\newtheorem{lemma}[theorem]{Lemma}
\newtheorem{proposition}[theorem]{Proposition}
\newtheorem{corollary}[theorem]{Corollary}
\newtheorem{conjecture}[theorem]{Conjecture}
\theoremstyle{definition}
\newtheorem{definition}[theorem]{Definition}
\newtheorem{remark}[theorem]{Remark}
\newtheorem{example}[theorem]{Example}

\title{Final pair identification: closing the support-size $2^a$ theorem\\
for $\lambda=(2^a,1^{n-2a})$}
\author{Clio}
\date{18 May 2026}

\begin{document}
\maketitle

\begin{abstract}
We close Conjecture~5.6 of \texttt{2026-05-17-trajectory-uniform.tex}
(``\textbf{May-17}''), the final residue of the support-size $2^a$
program for the family $\lambda=(2^a,1^{n-2a})$ at $n\ge 3a$.  The
proof tracks the column-$2$ entry set
$\sigma(P_k^+)$ through the trajectory and identifies it via
a closed-form description of $S_a^{(n)}$:
\[
   T_\varepsilon
   \;=\;
   \{i-1 : \varepsilon_i=L,\ 1\le i\le a\}
   \cup
   \{2a-i : \varepsilon_i=R,\ 1\le i\le a\}
   \qquad
   (\text{shifted by }-(n-2a+1)).
\]
The trajectory's evolution then breaks into three independent pieces:
the lower range (controlled by the May-17 $\sigma$-invariant), an
``upper'' range (controlled by a one-line recursion on $\tau$), and a
``middle'' pair $\{a-1,a\}$ in the shifted universe (controlled by
the case at the final step).  Matching the three pieces against the
closed form yields the final pair identification
\[
   \big\{\bar\sigma(P_{a-1}^+(T')),\,\bar\sigma(P_{a-1}^-(T'))\big\}
   \;=\;
   \{T_{\tau L},T_{\tau R}\},
   \qquad
   \tau:=\varepsilon(T')\in\{L,R\}^{a-1}.
\]

\textbf{Consequence.}  Combined with the May-17 dimension bound
$\dim B_{n-a+1}^{(\lambda)}V_\lambda\le 1$ (May-17~\S4.2) and the
generic non-vanishing argument, we obtain
\textbf{Conjecture~5.6 of May-17 unconditionally}:
$B_{n-a+1}^{(\lambda)}V_\lambda$ is exactly $1$-dimensional with
seminormal support $S_a^{(n)}$, $|S_a^{(n)}|=2^a$, all coefficients
nonzero in $\mathbb{Q}(q)$.  The rank-zero theorem for the family
$\lambda=(2^a,1^{n-2a})$ at every $a\ge 1$, $n\ge 3a$ holds (modulo
Phase~A endpoint).
\end{abstract}

\section{Setup and notation}\label{sec:setup}

We adopt the conventions of May-17.  Fix $a\ge 2$, $n\ge 3a$,
$\lambda=(2^a,1^{n-2a})\vdash n$, $\mu_1=(2^{a-1},1^{n-2a})\vdash n-2$.

\paragraph{Encoding.}  For $\sigma\in\binom{[2a]}{a}$ in the recursive
set $S_a$ (PROVE.md / May-14), the encoding
$\varepsilon:S_a\to\{L,R\}^a$ is defined recursively:
$\varepsilon_1=L$ if $0\in\sigma$, $R$ if $2a-1\in\sigma$; strip that
boundary element; shift the remainder down by $1$; recurse on
$S_{a-1}\subseteq[2a-2]$ to obtain $\varepsilon_2,\ldots,\varepsilon_a$.
For $T\in S_a^{(n)}$ (an SYT of $\lambda$ in the support, i.e.\ in
the $S_a$-indexed support stable at $n\ge 3a$), write
$\bar\sigma(T):=\sigma(T)-(n-2a+1)\subseteq[2a]$ (the column-$2$
entries shifted into the universe $\{0,1,\ldots,2a-1\}$).

For $T'\in S_{a-1}^{(n-2)}$, $\tau:=\varepsilon(T')\in\{L,R\}^{a-1}$.

\paragraph{Trajectory.}  As in May-17, the trajectory is
\[
   w_k(T')=(T_{n-1-k}{+}1)\cdots(T_{n-2}{+}1)\Phi_1(v_{T'}),
   \qquad k=0,1,\ldots,a-1,
\]
with $w_0=\Phi_1(v_{T'})$.  By Theorem~2.1 of May-17 (per-trajectory
invariant), $w_k$ has seminormal support $\{P_k^+,P_k^-\}$, a
non-degenerate $T_{n-1-k}$-$2$-block.  We use the May-17 convention
$P_0^+:=T_A(T')$ (the extension with $n-1$ at $(a,2)$, $n$ at
$(n-a,1)$), and $P_{k+1}^+:=$ the survivor of
$\{P_k^+,P_k^-\}$ under $(T_{n-2-k}{+}1)$ for $k\ge 0$.

We refer freely to the four-case structural Lemma~3.1 of May-17
(\emph{Kill/expand alternation}) with the case-labelling I, I.a, I.b,
II, II.a, II.b therein.

\section{Closed form for $S_a$}\label{sec:closed-form}

\begin{lemma}[Closed form]\label{lem:closed-form}
For every $a\ge 1$ and every $\varepsilon\in\{L,R\}^a$,
\[
   T_\varepsilon\;=\;
   \{i-1 : \varepsilon_i=L,\ 1\le i\le a\}
   \cup
   \{2a-i : \varepsilon_i=R,\ 1\le i\le a\}.
\]
In particular, $T_\varepsilon\cap\{0,1,\ldots,a-1\}=
\{i-1 : \varepsilon_i=L\}$ and
$T_\varepsilon\cap\{a,\ldots,2a-1\}=\{2a-i : \varepsilon_i=R\}$
(so the $L$- and $R$-bits ``fill'' disjoint halves of $[2a]$).
\end{lemma}

\begin{proof}
Induction on $a$.

\emph{Base $a=1$.}  $S_1=\{\{0\},\{1\}\}$.  $\varepsilon=L$ gives
$\{0\}=\{1-1\}$ (one $L$-contribution at $i=1$); $\varepsilon=R$
gives $\{1\}=\{2-1\}$ (one $R$-contribution at $i=1$).  $\checkmark$

\emph{Inductive step $a-1\to a$.}  Let
$\varepsilon=(\varepsilon_1,\ldots,\varepsilon_a)$.  By the $S_a$
recursion of PROVE.md / May-14, $T_\varepsilon=\{0\}\cup
\mathrm{shift}(T_{\varepsilon'})$ if $\varepsilon_1=L$ and
$T_\varepsilon=\mathrm{shift}(T_{\varepsilon'})\cup\{2a-1\}$ if
$\varepsilon_1=R$, where
$\varepsilon'=(\varepsilon_2,\ldots,\varepsilon_a)$ and
$T_{\varepsilon'}\in S_{a-1}$.

By the inductive hypothesis on $\varepsilon'\in\{L,R\}^{a-1}$,
\[
   T_{\varepsilon'}=
   \{i-1 : \varepsilon'_i=L\}\cup\{2(a-1)-i : \varepsilon'_i=R\}.
\]
Shifting by $+1$:
\[
   \mathrm{shift}(T_{\varepsilon'})
   =\{i : \varepsilon'_i=L\}\cup\{2a-1-i : \varepsilon'_i=R\}.
\]
Substituting $j=i+1$ (so $j$ ranges over $\{2,\ldots,a\}$ and
$\varepsilon'_i=\varepsilon_j$):
\[
   \mathrm{shift}(T_{\varepsilon'})
   =\{j-1 : \varepsilon_j=L,\ 2\le j\le a\}
   \cup
   \{2a-j : \varepsilon_j=R,\ 2\le j\le a\}.
\]
If $\varepsilon_1=L$, prepending $\{0\}=\{1-1 : \varepsilon_1=L\}$
extends the union to $j\in\{1,\ldots,a\}$, matching the claim.  If
$\varepsilon_1=R$, appending $\{2a-1\}=\{2a-1 : \varepsilon_1=R\}$
does the same.
\end{proof}

\begin{remark}
The closed form gives an elegant ``$L/R$-side'' partition of $[2a]$:
the $L$-bits contribute to the ``low half'' $\{0,\ldots,a-1\}$ with
the $i$-th $L$-bit at position $i-1$, and the $R$-bits contribute to
the ``high half'' $\{a,\ldots,2a-1\}$ with the $i$-th $R$-bit at
position $2a-i$ (mirror-symmetric).
\end{remark}

\section{Case alternation along the trajectory}\label{sec:case-alt}

We need to know, at each step $k$ of the trajectory, whether the
pair $\{P_k^+,P_k^-\}$ is in Case~I ($e_2$ in column $1$ of $P_k^+$,
$e_3$ in column $2$) or Case~II ($e_2$ in column $2$, $e_3$ in column
$1$), in the May-17 nomenclature.  The following determines the
case sequence from $\tau$.

\begin{lemma}[Case at step $k$]\label{lem:case}
For each $k\in\{0,1,\ldots,a-1\}$, the pair $\{P_k^+,P_k^-\}$ at the
start of step $k$ is:
\begin{itemize}[leftmargin=2em,topsep=0.3em,itemsep=0.2em]
\item $k=0$: Case~II.
\item $k\ge 1$: Case~I if $\tau_k=L$; Case~II if $\tau_k=R$.
\end{itemize}
\end{lemma}

\begin{proof}
\emph{Base $k=0$.}  $P_0^+=T_A(T')$ has $n-1$ at column~$2$ row $a$
and $n$ at column~$1$ row $n-a$.  So $e_2=n-1$ is in column~$2$ of
$P_0^+$ and $e_3=n$ is in column~$1$: Case~II.

\emph{Inductive step $k\to k+1$ (for $k\ge 0$).}  After the
kill/expand at step $k$ (May-17 Lemma~3.1), the new pair is
$\{P_{k+1}^+,P_{k+1}^-=s_{e_1^{(k)}}P_{k+1}^+\}$ where
$e_1^{(k)}=n-2-k$.  The new $e_2^{(k+1)}=n-1-(k+1)=n-2-k=e_1^{(k)}$.
So the case at step $k+1$ is determined by: is
$e_1^{(k)}\in\sigma(P_{k+1}^+)$?

By the May-17 Corollary on encoding-bit identification (Cor.~5.5):
$e_1^{(k)}=n-2-k\in\sigma(P_k^+)$ iff $\tau_{k+1}=R$.  The transposition
$s_{e_2^{(k)}}$ does not move $e_1^{(k)}$, so
$e_1^{(k)}\in\sigma(P_k^-)\iff e_1^{(k)}\in\sigma(P_k^+)$.  Hence in
either survivor (whether $P_{k+1}^+=P_k^+$ or $P_{k+1}^+=P_k^-$),
\[
   e_1^{(k)}\in\sigma(P_{k+1}^+)
   \iff
   \tau_{k+1}=R.
\]
The new $e_2^{(k+1)}=e_1^{(k)}$; ``$e_2^{(k+1)}\in\sigma$'' means
column~$2$, which is Case~II.  Hence Case~II at step $k+1$ iff
$\tau_{k+1}=R$, i.e.\ Case~I iff $\tau_{k+1}=L$.
\end{proof}

\begin{remark}
The case is governed by the sub-case (a/b) of the \emph{previous}
step: sub-case~(a) $\Rightarrow$ Case~I next; sub-case~(b)
$\Rightarrow$ Case~II next.  Equivalently, sub-case~(a) and Case~I
are both ``column~$1$'' phenomena, while sub-case~(b) and Case~II
are both ``column~$2$''.
\end{remark}

\section{Three-piece $\sigma$-decomposition}\label{sec:decomp}

For each $k\in\{0,1,\ldots,a-1\}$, partition
\[
   \{n-2a+1,\ldots,n\}\;=\;
   \underbrace{\{n-2a+1,\ldots,n-2-k\}}_{\text{lower}}
   \;\sqcup\;
   \underbrace{\{n-1-k,n-k\}}_{\text{middle}}
   \;\sqcup\;
   \underbrace{\{n-k+1,\ldots,n\}}_{\text{upper}}.
\]
For $k=a-1$ the three blocks have sizes $a-1$, $2$, $a-1$.  Set:
\begin{align}
   A_k & := \sigma(P_k^+)\cap\{n-2a+1,\ldots,n-2-k\}, \\
   M_k & := \sigma(P_k^+)\cap\{n-1-k,n-k\}, \\
   U_k & := \sigma(P_k^+)\cap\{n-k+1,\ldots,n\}.
\end{align}
Then $\sigma(P_k^+)=A_k\sqcup M_k\sqcup U_k$ (and similarly with
$P_k^-$, with $M_k$ swapped and $A_k,U_k$ unchanged since
$P_k^-=s_{n-1-k}P_k^+$ only moves entries in $\{n-1-k,n-k\}$).

\begin{lemma}[Three-piece formula]\label{lem:three-piece}
For each $k\in\{0,1,\ldots,a-1\}$:
\begin{enumerate}[leftmargin=2em,topsep=0.3em,itemsep=0.2em,label=(\roman*)]
\item $A_k=\sigma(T')\cap\{n-2a+1,\ldots,n-2-k\}$.  (Lower part.)
\item $U_k=\{n-j : 0\le j\le k-1,\ \tau_{j+1}=R\}$.  (Upper part.)
\item $M_k=\{n-1-k\}$ if step $k$ is Case~II, and $\{n-k\}$ if Case~I.
\end{enumerate}
\end{lemma}

\begin{proof}
\textbf{(i)}  This is the May-17 $\sigma$-invariant (Proposition~5.1).

\medskip\noindent\textbf{(iii)}  Immediate from the definitions of
Cases~I, II in May-17 Lemma~3.1: at the start of step $k$ the pair is
a $T_{n-1-k}$-$2$-block, $P_k^+$ has $e_2=n-1-k$ and $e_3=n-k$ at
different columns of $\lambda$, with one in column~$2$ (yielding the
$\sigma$-membership) and the other in column~$1$.

\medskip\noindent\textbf{(ii)}  Induction on $k$.  Base $k=0$:
$U_0=\sigma(P_0^+)\cap\{n+1,\ldots,n\}=\emptyset$, matching the
empty product on the right.

\emph{Inductive step $k\to k+1$.}  $U_{k+1}\cap\{n-k+1,\ldots,n\}
=\sigma(P_{k+1}^+)\cap\{n-k+1,\ldots,n\}$.  Now $P_{k+1}^+$ is the
survivor of $\{P_k^+,P_k^-\}$, both of which agree on entries
\emph{outside} $\{e_2^{(k)},e_3^{(k)}\}=\{n-1-k,n-k\}$, in particular
on $\{n-k+1,\ldots,n\}$.  So
\[
   \sigma(P_{k+1}^+)\cap\{n-k+1,\ldots,n\}
   =\sigma(P_k^+)\cap\{n-k+1,\ldots,n\}
   =U_k.
\]
Then $U_{k+1}=U_k\cup\big(\{n-k\}\cap\sigma(P_{k+1}^+)\big)$
(adding membership at the new ``boundary'' position $n-k$, which
moves from middle at step $k$ to upper at step $k+1$).

\emph{Whether $n-k\in\sigma(P_{k+1}^+)$.}  This depends on the
sub-case at step $k$.  By May-17 Lemma~3.1 case analysis:
\begin{itemize}[leftmargin=2em,topsep=0.3em,itemsep=0.2em]
\item Case~I.a: $P_{k+1}^+=P_k^-$.  In Case~I, $e_3=n-k$ is in
column~$2$ of $P_k^+$, hence in column~$1$ of $P_k^-=s_{e_2}P_k^+$.
So $n-k\notin\sigma(P_{k+1}^+)$.
\item Case~I.b: $P_{k+1}^+=P_k^+$, $e_3$ in column~$2$.  So
$n-k\in\sigma(P_{k+1}^+)$.
\item Case~II.a: $P_{k+1}^+=P_k^+$.  In Case~II, $e_3$ in column~$1$.
So $n-k\notin\sigma(P_{k+1}^+)$.
\item Case~II.b: $P_{k+1}^+=P_k^-$, $e_3$ in column~$2$ of $P_k^-$.
So $n-k\in\sigma(P_{k+1}^+)$.
\end{itemize}
In all four: $n-k\in\sigma(P_{k+1}^+)$ iff the sub-case is~(b),
iff (by May-17 Cor.~5.5) $\tau_{k+1}=R$.

\medskip
Hence $U_{k+1}=U_k\cup(\{n-k\}$ if $\tau_{k+1}=R$, else $\emptyset)$.
Combined with the inductive hypothesis,
\[
   U_{k+1}=\{n-j : 0\le j\le k,\ \tau_{j+1}=R\}.
   \qedhere
\]
\end{proof}

\section{Final pair identification}\label{sec:main}

\begin{theorem}[Final pair $=\{T_{\tau L},T_{\tau R}\}$]\label{thm:final-pair}
For every $a\ge 2$, every $n\ge 3a$, every $T'\in S_{a-1}^{(n-2)}$
with encoding $\tau=\varepsilon(T')\in\{L,R\}^{a-1}$,
\[
   \big\{\bar\sigma(P_{a-1}^+(T')),\,\bar\sigma(P_{a-1}^-(T'))\big\}
   \;=\;\{T_{\tau L},T_{\tau R}\}.
\]
In particular, $\bar\sigma(P_{a-1}^+),\bar\sigma(P_{a-1}^-)\in S_a$
with encoding prefix $\tau$ and last bit varying over $\{L,R\}$.
\end{theorem}

\begin{proof}
Apply Lemma~\ref{lem:three-piece} at $k=a-1$ and shift by
$-(n-2a+1)$:
\begin{align*}
   \bar A_{a-1} & =\bar\sigma(T')\cap\{0,1,\ldots,a-2\}, \\
   \bar M_{a-1} & =\begin{cases}\{a-1\} & \text{Case~II at step }a-1\\
                                \{a\}   & \text{Case~I at step }a-1\end{cases}, \\
   \bar U_{a-1} & =\{2a-1-j : 0\le j\le a-2,\ \tau_{j+1}=R\}.
\end{align*}
By Lemma~\ref{lem:case}, the case at step $a-1$ is:
Case~I iff $\tau_{a-1}=L$ (so $\bar M_{a-1}=\{a\}$ in that case),
Case~II iff $\tau_{a-1}=R$ (so $\bar M_{a-1}=\{a-1\}$).

\paragraph{Simplifying $\bar A_{a-1}$ and $\bar U_{a-1}$.}
By the closed-form Lemma~\ref{lem:closed-form} applied to
$\bar\sigma(T')=T_\tau\in S_{a-1}$,
\[
   \bar\sigma(T')
   =\{i-1 : \tau_i=L\}\cup\{2(a-1)-i : \tau_i=R\}
   \quad(1\le i\le a-1).
\]
The $L$-contributions lie in $\{0,\ldots,a-2\}$ and the
$R$-contributions in $\{a-1,\ldots,2a-3\}$
(Lemma~\ref{lem:closed-form} Remark).  Restricting to
$\{0,\ldots,a-2\}$:
\[
   \bar A_{a-1}=\{i-1 : \tau_i=L,\ 1\le i\le a-1\}.
\]
For $\bar U_{a-1}$, substitute $i=j+1$:
\[
   \bar U_{a-1}=\{2a-i : \tau_i=R,\ 1\le i\le a-1\}.
\]

\paragraph{Comparison with the closed form for $S_a$.}  By
Lemma~\ref{lem:closed-form} applied to $\varepsilon=\tau\cdot b\in
\{L,R\}^a$ where $b\in\{L,R\}$ is the final bit:
\begin{align*}
   T_{\tau L} & =\{i-1 : \tau_i=L,\ 1\le i\le a-1\}\cup\{a-1\}
   \cup\{2a-i : \tau_i=R,\ 1\le i\le a-1\}, \\
   T_{\tau R} & =\{i-1 : \tau_i=L,\ 1\le i\le a-1\}\cup\{a\}
   \cup\{2a-i : \tau_i=R,\ 1\le i\le a-1\}
\end{align*}
(the bit at position $a$ contributes $a-1$ if $L$, $a$ if $R$).
Equivalently,
\[
   T_{\tau L}=\bar A_{a-1}\cup\{a-1\}\cup\bar U_{a-1},\qquad
   T_{\tau R}=\bar A_{a-1}\cup\{a\}\cup\bar U_{a-1}.
\]

\paragraph{Two cases.}
\begin{itemize}[leftmargin=2em,topsep=0.3em,itemsep=0.2em]
\item $\tau_{a-1}=L$: Case~I at step $a-1$.
  $\bar\sigma(P_{a-1}^+)=\bar A_{a-1}\cup\{a\}\cup\bar U_{a-1}=T_{\tau R}$.
  And $\bar\sigma(P_{a-1}^-)=\bar A_{a-1}\cup\{a-1\}\cup\bar U_{a-1}=T_{\tau L}$
  (swap the $\bar M$ piece, $\bar A$ and $\bar U$ unchanged).
\item $\tau_{a-1}=R$: Case~II at step $a-1$.
  $\bar\sigma(P_{a-1}^+)=T_{\tau L}$, $\bar\sigma(P_{a-1}^-)=T_{\tau R}$.
\end{itemize}
Either way, $\{\bar\sigma(P_{a-1}^+),\bar\sigma(P_{a-1}^-)\}
=\{T_{\tau L},T_{\tau R}\}$.
\end{proof}

\begin{remark}[Which-is-which depends on the path]
The assignment $P_{a-1}^+\leftrightarrow T_{\tau L}$ vs.\ $T_{\tau R}$
depends on $\tau_{a-1}$.  At $a=2$ with $\tau=L$ (so $\tau_{a-1}=L$):
$P_1^+\leftrightarrow T_{LR}$ (verified in May-17 \S5.2 by hand).
At $a=2,\tau=R$: $P_1^+\leftrightarrow T_{RL}$.  The asymmetry comes
from the case-flip rule: $\tau_{a-1}=L$ forces Case~I at the last
step, where $P^+$ ends with $\bar M=\{a\}$ (matching $T_{\tau R}$).
\end{remark}

\section{Consequences}\label{sec:consequences}

\begin{theorem}[Support $=S_a^{(n)}$ unconditionally]\label{thm:support}
For every $a\ge 2$ and $n\ge 3a$,
$B_{n-a+1}^{(\lambda)}V_\lambda$ is exactly $1$-dimensional with
seminormal support $S_a^{(n)}$ (of size $|S_a^{(n)}|=2^a$),
all coefficients in $\mathbb{Q}(q)$ nonzero.
\end{theorem}

\begin{proof}
The dimension bound $\dim B_{n-a+1}^{(\lambda)}V_\lambda\le 1$ was
established in May-17 \S4.2.  By May-15-gen Theorem~3.2,
\[
   B_{n-a+1}^{(\lambda)}V_\lambda
   =\mathbb{Q}(q)\cdot\sum_{T'\in\mathrm{supp}(u^{(\mu_1)})}c_{T'}\,w_{a-1}(T'),
\]
where $u^{(\mu_1)}$ generates $B_{n-a}^{(\mu_1)}V_{\mu_1}$, which by
induction at $a-1$ has support exactly $S_{a-1}^{(n-2)}$ and is
$1$-dimensional, so $|\mathrm{supp}(u^{(\mu_1)})|=2^{a-1}$ with
nonzero coefficients $c_{T'}$.

By Theorem~\ref{thm:final-pair}, $w_{a-1}(T')$ has support
$\{T_{\tau L},T_{\tau R}\}$ where $\tau=\varepsilon(T')$.  As $T'$
ranges over $S_{a-1}^{(n-2)}$, $\tau$ ranges over $\{L,R\}^{a-1}$
bijectively, hence
\[
   \bigsqcup_{T'\in S_{a-1}^{(n-2)}}\{T_{\tau(T')L},T_{\tau(T')R}\}
   \;=\;S_a^{(n)},
\]
because $\varepsilon:S_a\to\{L,R\}^a$ is a bijection and the pairs
$\{T_{\tau L},T_{\tau R}\}$ partition $S_a$ according to their first
$a-1$ encoding bits.  The supports are pairwise disjoint (different
$\tau$'s give different prefixes, hence different pairs).
Therefore
$\bigcup_{T'}\mathrm{supp}(w_{a-1}(T'))=S_a^{(n)}$.

\emph{Nonzero coefficients.}  By May-17 (P4) at $k=a-1$, both
seminormal coefficients of $w_{a-1}(T')$ are nonzero in
$\mathbb{Q}(q)$.  By the bijection $T'\leftrightarrow\tau$ and the
pair-partition, the linear combination
$\sum_{T'}c_{T'}w_{a-1}(T')$ has coefficient at each
$T_\varepsilon\in S_a$ equal to a single (nonzero) contribution from
the unique $T'$ with $\varepsilon(T')$-matching prefix.  Hence the
support is exactly $S_a^{(n)}$ with all $\mathbb{Q}(q)$-coefficients
nonzero, and the sum is nonzero, confirming dimension $1$.
\end{proof}

\begin{corollary}[Rank-zero theorem at $(2^a,1^{n-2a})$, unconditional
on conjectural support]\label{cor:rank-zero}
For every $a\ge 1$ and $n\ge 3a$,
$\Pi^{S_n}|_{V_{(2^a,1^{n-2a})}}=0$, conditional only on the Phase~A
endpoint $R'_nV_\lambda=E_{n-1}^+|_{V_\lambda}$ at general $a$.
\end{corollary}

\begin{proof}
May-17 Theorem~4.1 gives this conditionally on Conjecture~5.6.
Theorem~\ref{thm:support} above closes that conjecture.
\end{proof}

\section{Computational verification}\label{sec:verify}

We verified the final pair identification by direct $\sigma$-level
trajectory simulation: for each $T'\in S_{a-1}^{(n-2)}$, simulate the
kill/expand trajectory using only the structural rule (kill iff $e_1$
and $e_2$ in the same column), and compare the final pair against the
predicted $\{T_{\tau L},T_{\tau R}\}$ from the closed form.

\begin{center}
\begin{tabular}{c c c l}
\toprule
$a$ & $n$ & $|S_{a-1}^{(n-2)}|$ & result \\
\midrule
2 & 6, 7, 8       & 2 each   & PASS (all $2$ trajectories match) \\
3 & 9, 10, 11     & 4 each   & PASS (all $4$ trajectories match) \\
4 & 12, 13, 14    & 8 each   & PASS (all $8$ trajectories match) \\
5 & 15, 16        & 16 each  & PASS (all $16$ trajectories match) \\
6 & 18            & 32       & PASS (all $32$ trajectories match) \\
\bottomrule
\end{tabular}
\end{center}

Script: \texttt{\textasciitilde/projects/scratch/2026-05-18-final-pair/verify\_trajectory.py}.
Total of $12$ cases at $a\in\{2,3,4,5,6\}$, covering $2+2+2+4+4+4+8+8+8+16+16+32=106$
trajectory simulations.  All match the closed-form prediction
$\{T_{\tau L},T_{\tau R}\}$.

A separate script
(\texttt{verify\_closed\_form.py}) verified the closed-form $S_a$
formula against the recursive definition for $a\le 6$ (all $|S_a|=2^a$
elements, with encoding round-trips through both directions).

We also performed hand verifications at $a=3$ in
\S\ref{sec:examples}.

\section{Examples}\label{sec:examples}

\subsection{$a=3$, $n=9$, $\tau=LL$}

$T'$ with $\bar\sigma(T')=\{0,1\}$, i.e.\ $\sigma(T')=\{4,5\}$.  The
trajectory:
\begin{enumerate}[leftmargin=2em,topsep=0.3em,itemsep=0.2em]
\item $P_0^+=T_A$, $\sigma=\{4,5,8\}$.  Step~$0$ ($e_1=7$): $7\notin
\sigma(T')$, sub-case~(a), Case~II.a.  Survivor $P_1^+=T_A$.
\item $P_1^+=T_A$.  Step~$1$ ($e_1=6$): $6\notin\sigma(T_A)$,
sub-case~(a), Case~I.a (Case~I per Lemma~\ref{lem:case} since
$\tau_1=L$).  Survivor $P_2^+=P_1^-=s_7T_A$, $\sigma=\{4,5,7\}$.
\item Final pair: $\{s_7T_A,s_6s_7T_A\}$, $\sigma$-values
$\{\{4,5,7\},\{4,5,6\}\}$.  Shifted by $-4$:
$\{\{0,1,3\},\{0,1,2\}\}=\{T_{LLR},T_{LLL}\}=\{T_{\tau R},T_{\tau L\}}$.
$\checkmark$
\end{enumerate}

\subsection{$a=3$, $n=9$, $\tau=LR$}

$T'$ with $\bar\sigma(T')=\{0,2\}$, i.e.\ $\sigma(T')=\{4,6\}$.
\begin{enumerate}[leftmargin=2em,topsep=0.3em,itemsep=0.2em]
\item $P_0^+=T_A$, $\sigma=\{4,6,8\}$.  Step~$0$ ($e_1=7$): $7\notin
\sigma(T')$, Case~II.a.  Survivor $P_1^+=T_A$.
\item $P_1^+=T_A$.  Step~$1$ ($e_1=6$): $6\in\sigma(T_A)$, sub-case~(b),
Case~I.b.  Survivor $P_2^+=P_1^+=T_A$, $\sigma=\{4,6,8\}$.
\item Final pair: $\{T_A,s_6T_A\}$, $\sigma$-values
$\{\{4,6,8\},\{4,7,8\}\}$.  Shifted:
$\{\{0,2,4\},\{0,3,4\}\}=\{T_{LRL},T_{LRR}\}=\{T_{\tau L},T_{\tau R}\}$.
$\checkmark$
\end{enumerate}

\subsection{Formula check}

By Theorem~\ref{thm:final-pair}, with $\tau=LR$:
\begin{align*}
   \bar A_{a-1} & = \{i-1 : \tau_i=L,\ 1\le i\le 2\} = \{0\}, \\
   \bar U_{a-1} & = \{2a-i : \tau_i=R,\ 1\le i\le 2\} = \{4\}, \\
   \bar M_{a-1}(P^+) & = \{a-1\}=\{2\}\quad\text{(Case II, since }\tau_{a-1}=\tau_2=R\text{).}
\end{align*}
Predicting $\bar\sigma(P_2^+)=\{0,2,4\}=T_{\tau L}=T_{LRL}$. $\checkmark$

\section{Honest assessment}\label{sec:honest}

\paragraph{What this paper rigorously proves.}
\begin{enumerate}[leftmargin=2em,topsep=0.3em,itemsep=0.2em]
\item Closed form (Lemma~\ref{lem:closed-form}) for elements of
$S_a$ in terms of their encoding.  Elementary induction.
\item Case alternation at each step (Lemma~\ref{lem:case}): each
step's case (I vs.\ II) is determined by the corresponding encoding
bit, with step $0$ always Case~II.  Follows from May-17 Cor.~5.5.
\item Three-piece $\sigma$-decomposition (Lemma~\ref{lem:three-piece}):
$\sigma(P_k^+)$ splits into lower (May-17 invariant), middle (case),
upper ($\tau$-tracking).
\item Final pair identification (Theorem~\ref{thm:final-pair}):
$\{\bar\sigma(P_{a-1}^\pm)\}=\{T_{\tau L},T_{\tau R}\}$.
\item Support identification (Theorem~\ref{thm:support}): the
support of $B_{n-a+1}^{(\lambda)}V_\lambda$ is exactly $S_a^{(n)}$,
$|S_a^{(n)}|=2^a$, dimension $1$, all coefficients nonzero.
\item Rank-zero theorem (Corollary~\ref{cor:rank-zero}) for the
family $\lambda=(2^a,1^{n-2a})$ at every $a\ge 1$, $n\ge 3a$.
Conditional on Phase~A endpoint.
\end{enumerate}

\paragraph{What remains conditional.}
\begin{enumerate}[leftmargin=2em,topsep=0.3em,itemsep=0.2em]
\item Phase~A endpoint at general $a\ge 4$: the assertion
$R'_nV_\lambda=E_{n-1}^+|_{V_\lambda}$.  Proved at $a=2$ (May-13
Prop.~3.1) and $a=3$ (May-14 Prop.~3.1); asserted mechanical for
$a\ge 4$ in May-15-gen Theorem~2.1.  Computationally verified for
$a=4,5$ at small $n$.  This is an orthogonal piece of work,
structurally independent of the per-trajectory framework.
\item Generic non-vanishing of the summed image
$\sum_{T'}c_{T'}w_{a-1}(T')$: handled in Theorem~\ref{thm:support}
via the disjoint-pair argument and (P4) of the May-17 invariant.
\end{enumerate}

\paragraph{What's NOT conditional anymore.}  The support-size $2^a$
conjecture and the L/R encoding identification — the explicit
combinatorial content of the May-13/May-15/May-16 case enumerations
— is now proved \emph{uniformly} via:
\begin{itemize}[leftmargin=2em,topsep=0.3em,itemsep=0.2em]
\item The May-17 four-case structural lemma (per-trajectory invariant).
\item The May-17 $\sigma$-invariant (lower-range freezing).
\item The closed form for $S_a$ (Lemma~\ref{lem:closed-form}).
\item The three-piece $\sigma$-decomposition (Lemma~\ref{lem:three-piece}).
\item Theorem~\ref{thm:final-pair}.
\end{itemize}
No case-by-case enumeration; no encoding-bit bookkeeping beyond the
clean ``$\bar M_{a-1}\in\{\{a-1\},\{a\}\}$'' final-step distinction.

\paragraph{Big picture.}  The May-13/15/16/17 program has been
moving toward exactly this clean statement.  May-13 enumerated by
hand at $a=2$.  May-15-a3 enumerated by hand at $a=3$.  May-16
identified the per-trajectory framework and conjectured the uniform
support flow.  May-17 closed the uniform per-trajectory invariant
and the $\sigma$-invariant, isolating the residue.  This paper
closes the residue via the closed form for $S_a$ and the upper-$U_k$
tracking.

The proof is short because the underlying object is well-structured:
elements of $S_a$ have a clean closed form, and $\sigma(P_k^+)$
decomposes into three independent pieces that each have a clean
description in terms of $\tau$.  The original case enumeration was
illusorily complex; once the right invariants are isolated, the
identification is mechanical.

\end{document}
